%------------------------------------------------------------------------------%
\section{Funções Trigonométricas Inversas}

As funções seno, cosseno e tangente não são bijetoras, mas uma restrição no
domínio das mesmas nos possibilita definir suas inversas. Essa escolha de
domínio não é única, mas adotaremos as escolhas padrões, que são os intervalos
que contém os ângulos notáveis e cujas imagens das funções restritas são iguais
às imagens das funções sem restrição. Os gráficos das funções inversas são
obtidos a partir dos gráficos das funções originais através da reflexão do
gráfico da função original em torno da reta $y=x$. Por isso apresentaremos três
imagens: o gráfico das funções trigonométricas restritas, esse gráfico
juntamente com sua reflexão e o gráfico das funções trigonométricas inversas.

%------------------------------------------------------------------------------%
\subsection{Arco Seno}

A função
\[
  \sin\colon \left[-\frac{\pi}{2}, \frac{\pi}{2}\right] \to [-1, 1]
\]
é bijetora e sua inversa é denominada arco seno
\[
  \arcsin\colon [-1, 1] \to \left[-\frac{\pi}{2}, \frac{\pi}{2}\right].
\]
Por definição,
\[
  \arcsin x=y \qquad\Leftrightarrow\qquad \sin y=x.
\]

Por serem inversas uma da outra, essas duas funções satisfazem:
\begin{gather*}
  \sin(\arcsin(x))=x, \quad\forall\quad x~\in [-1,1]. \\[2mm]
  \arcsin(\sin(x))=x, \quad\forall\quad x~\in \left[-\frac{\pi}{2}, \frac{\pi}{2}\right].
\end{gather*}

Para encontrarmos o gráfico da função arcoseno, vamos analisar o  gráfico da
função seno em \(\left[-\sfrac{\pi}{2}, \sfrac{\pi}{2}\right]\) como dado a seguir:

\begin{figure}
  \centering
  \includegraphics[width=0.6\textwidth,trim={1cm 10cm 12cm 3cm},clip]{sen1}
  \caption{???}
\end{figure}

A reflexão do gráfico de seno em torno no eixo $x$:

\begin{figure}
  \centering
  \includegraphics[width=0.6\textwidth,trim={1cm 4cm 12cm 2cm},clip]{sen2}
  \caption{???}
\end{figure}

E finalmente, o gráfico do arco seno:

\begin{figure}
  \centering
  \includegraphics[width=0.6\textwidth,trim={1cm 6cm 12cm 3cm},clip]{sen4}
  \caption{???}
\end{figure}

\begin{example}
  Calcule o valor de $\arcsin x$ para $x$ assumindo os seguintes valores:
  \[
    -\frac{\sqrt{3}}{2}, \quad
    0,                   \quad
    \frac{\sqrt{2}}{2},  \quad
    \frac{1}{2},         \quad
    1
  \]

  \solution
  Temos que:
  \begin{gather*}
    \arcsin\left(-\frac{\sqrt{3}}{2}\right) = -\frac{\pi}{3} \\
    \arcsin(0)                              = 0              \\
    \arcsin\left(\frac{\sqrt{2}}{2}\right)  = \frac{\pi}{4}  \\
    \arcsin\left(\frac{1}{2}\right)         = \frac{\pi}{6}  \\
    \arcsin(1)                              = \frac{\pi}{2}
  \end{gather*}
  Vale lembrar que o valor de $\arcsin(x)$ está restrito ao intervalo
  \[
    \left[-\frac{\pi}{2}, \frac{\pi}{2}\right].
  \]
  Por exemplo, no caso de
  \[
    \arcsin\left(-\frac{\sqrt{3}}{2}\right),
  \]
  o valor $-\frac{\pi}{3}$ está
  nesse intervalo, mas o valor $\frac{2\pi}{3}$ também é uma solução da equação
  \[
    \sin(x)=-\frac{\sqrt{3}}{2},
  \]
  porém não está no intervalo
  \[
    \left[-\frac{\pi}{2},\frac{\pi}{2}\right].
  \]
  Por isso, é importante sempre verificar se a solução encontrada está
  dentro do intervalo de definição de $\arcsin(x)$.
\end{example}

%------------------------------------------------------------------------------%
\subsection{Arco Cosseno}

A função
\[
  \cos\colon [0, \pi] \to [-1, 1]
\]
é bijetora e sua inversa é denominada arco cosseno:
\[
  \arccos\colon [-1, 1] \to [0, \pi].
\]
Por definição,
\[
  \arccos x = y \Leftrightarrow \cos y = x.
\]

Por serem inversas uma da outra, essas duas funções satisfazem
\begin{gather*}
  \cos(\arccos(x))=x, \quad\forall\quad x~\in [-1, 1]. \\
  \arccos(\cos(x))=x, \quad\forall\quad x~\in [0, \pi].
\end{gather*}

A seguir o gráfico da função cosseno em $[0, \pi]$
e de sua reflexão em torno de $y=x$.

\begin{figure}
  \centering
  \includegraphics[width=0.4\textwidth,trim={1cm 4cm 12cm 2cm},clip]{cos2}
  \qquad
  \includegraphics[width=0.4\textwidth,trim={1cm 3cm 12cm 1cm},clip]{cos3}
  \caption{???}
\end{figure}

O gráfico do arco cosseno.

\begin{figure}
  \centering
  \includegraphics[height=60mm]{cos4}
  \caption{???}
\end{figure}

\begin{example}
  Calcule o valor de $\arccos x$ para $x$ assumindo os seguintes valores:
  \[
    -\frac{\sqrt{3}}{2}, \quad
    0,                   \quad
    \frac{\sqrt{2}}{2},  \quad
    \frac{1}{2},         \quad
    1
  \]

  \solution
  Temos que:
  \begin{gather*}
    \arccos\left(-\frac{\sqrt{3}}{2}\right)=\frac{5\pi}{6}. \\
    \arccos(0)=\frac{\pi}{2}.                               \\
    \arccos\left(\frac{\sqrt{2}}{2}\right)=\frac{\pi}{4}.   \\
    \arccos\left(\frac{1}{2}\right)=\frac{\pi}{3}.          \\
    \arccos(1)=0.
  \end{gather*}
  Vale lembrar que o valor de $\arccos(x)$ está restrito ao intervalo $[0, \pi]$.
  Por exemplo, no caso de
  \[
    \arccos\left(-\frac{\sqrt{3}}{2}\right),
  \]
  o valor
  $\sfrac{7\pi}{6}$ é uma solução da equação
  \[
    \cos(x)=-\frac{\sqrt{3}}{2},
  \]
  porém não está no intervalo $[0, \pi]$.
  Por isso, é importante sempre verificar se a
  solução encontrada está dentro do intervalo de definição de $\arccos(x)$.
\end{example}

%------------------------------------------------------------------------------%
\subsection{Arco Tangente}

A função
\[
  \tan\colon (-\pi/2, \pi/2) \to \R
\]
é bijetora e sua inversa é denominada arco tangente
\[
  \arctan\colon \R \to (-\pi/2, \pi/2).
\]
Por definição,
\[
  \arctan x=y \Leftrightarrow \tan y=x.
\]

Por serem inversas uma da outra, essas duas funções satisfazem:
\begin{align*}
  \tan(\arctan(x))=x, \quad\forall\quad x~\in \R. \\
  \arctan(\tan(x))=x, \quad\forall\quad x~\in \left(-\frac{\pi}{2}, \frac{\pi}{2}\right).
\end{align*}

A seguir o gráfico da função tangente em
\(\left(-\sfrac{\pi}{2}, \sfrac{\pi}{2}\right)\)
e de sua reflexão em torno de $y=x$

\begin{figure}
  \centering
  \includegraphics[width=0.4\textwidth,trim={1cm 2cm 12cm 2cm},clip]{tg1} \qquad
  \includegraphics[width=0.4\textwidth,trim={1cm 1cm 10cm 1cm},clip]{tg2}
  \caption{???}
\end{figure}

O gráfico do arco tangente.

\begin{figure}
  \centering
  \includegraphics[width=0.6\textwidth,trim={1cm 6cm 10cm 3cm},clip]{tg4}
  \caption{???}
\end{figure}

\begin{example}
  Calcule o valor de $\arctan x$ para $x$ assumindo os seguintes valores:
  \[
    -\frac{\sqrt{3}}{3}, \quad
    0,                   \quad
    \sqrt{3},            \quad
    1
  \]

  \solution
  Temos que:
  \begin{gather*}
    \arctan\left(-\frac{\sqrt{3}}{3}\right)=-\frac{\pi}{6}. \\
    \arctan\left(0\right)=0.                                \\
    \arctan\left(\sqrt{3}\right)= \frac{\pi}{3}.            \\
    \arctan\left(1\right)=\frac{\pi}{4}.
  \end{gather*}
  Vale lembrar que o valor de $\arctan(x)$ está restrito ao intervalo
  \[
    \left(-\frac{\pi}{2},\frac{\pi}{2}\right).
  \]
  Por exemplo, no caso de
  \[
    \arctan\left(-\frac{\sqrt{3}}{3}\right),
  \]
  a solução encontrada
  \[
    -\frac{\pi}{6}
  \]
  está dentro do intervalo
  \[
    \left(-\frac{\pi}{2},\frac{\pi}{2}\right),
  \]
  porém é importante sempre verificar se a solução encontrada está dentro
  do intervalo de definição de $\arctan(x)$.
\end{example}

%------------------------------------------------------------------------------%
